% GENERATED FILE. Edit canonical lessons, then run npm run generate:books.
% canonical-source-hash: fnv1a64:5b489302261dbcc5
% canonical-lessons: KA-C139-kamala, KA-C139-godhi, KA-C139-doni, KA-C139-kanniru, KA-C139-mise

\chapter{Lotus, Wheat, Boat, Tears, Moustache}
\label{ch:ka-lotus-wheat-boat-tears-moustache}
\hlchaptermodality{139}

\begin{chapteropening}
\textbf{By the end of this chapter:} \emph{I can say a lotus, wheat, a boat, tears and a moustache.}

Complete the last lesson of chapter 139: i can say a lotus, wheat, a boat, tears and a moustache.
\end{chapteropening}

\section[kamala]{\kn{ಕಮಲ} (kamala) — a lotus}

\label{lesson:KA-C139-kamala}



\begin{quote}
Before the new one: say the Kannada for socks, then the Kannada for a coat.
\end{quote}

\subsection*{You'll want to know: \kn{ಕಮಲ}}
\textbf{\kn{ಕಮಲ}} — \emph{kamala} — \textquotedblleft{}a lotus\textquotedblright{}.

\begin{grammarlens}[title={What you've built}]
One more word for this chapter.
\end{grammarlens}

\subsection*{Your turn}
\begin{itemize}
\raggedright
  \item \emph{Say it:} \emph{kamala}
  \item \emph{Say it:} \emph{kamala}, once more
  \item \emph{Say it:} say \emph{kamala}
  \item \emph{Recall:} say the Kannada for socks, then the Kannada for a coat, then say \emph{kamala} again
\end{itemize}

\begin{tcolorbox}[breakable,colback=teal!4,colframe=teal!35!black,title={Before you move on}]
What does \kn{ಕಮಲ} mean? (\textquotedblleft{}A lotus\textquotedblright{}.) Say it once more.
\end{tcolorbox}
\section[gōdhi]{\kn{ಗೋಧಿ} (gōdhi) — wheat}

\label{lesson:KA-C139-godhi}



\begin{quote}
Before the new one: say the Kannada for a coat, then the Kannada for a lotus.
\end{quote}

\subsection*{You'll want to know: \kn{ಗೋಧಿ}}
\textbf{\kn{ಗೋಧಿ}} — \emph{gōdhi} — \textquotedblleft{}wheat\textquotedblright{}.

\begin{grammarlens}[title={What you've built}]
One more word for this chapter.
\end{grammarlens}

\subsection*{Your turn}
\begin{itemize}
\raggedright
  \item \emph{Say it:} \emph{gōdhi}
  \item \emph{Say it:} \emph{gōdhi}, once more
  \item \emph{Say it:} say \emph{gōdhi}
  \item \emph{Recall:} say the Kannada for a coat, then the Kannada for a lotus, then say \emph{gōdhi} again
\end{itemize}

\begin{tcolorbox}[breakable,colback=teal!4,colframe=teal!35!black,title={Before you move on}]
What does \kn{ಗೋಧಿ} mean? (\textquotedblleft{}Wheat\textquotedblright{}.) Say it once more.
\end{tcolorbox}
\section[dōṇi]{\kn{ದೋಣಿ} (dōṇi) — a boat}

\label{lesson:KA-C139-doni}



\begin{quote}
Before the new one: say the Kannada for a lotus, then the Kannada for wheat.
\end{quote}

\subsection*{You'll want to know: \kn{ದೋಣಿ}}
\textbf{\kn{ದೋಣಿ}} — \emph{dōṇi} — \textquotedblleft{}a boat\textquotedblright{}.

\begin{grammarlens}[title={What you've built}]
One more word for this chapter.
\end{grammarlens}

\subsection*{Your turn}
\begin{itemize}
\raggedright
  \item \emph{Say it:} \emph{dōṇi}
  \item \emph{Say it:} \emph{dōṇi}, once more
  \item \emph{Say it:} say \emph{dōṇi}
  \item \emph{Recall:} say the Kannada for a lotus, then the Kannada for wheat, then say \emph{dōṇi} again
\end{itemize}

\begin{tcolorbox}[breakable,colback=teal!4,colframe=teal!35!black,title={Before you move on}]
What does \kn{ದೋಣಿ} mean? (\textquotedblleft{}A boat\textquotedblright{}.) Say it once more.
\end{tcolorbox}
\section[kaṇṇīru]{\kn{ಕಣ್ಣೀರು} (kaṇṇīru) — tears}

\label{lesson:KA-C139-kanniru}



\begin{quote}
Before the new one: say the Kannada for wheat, then the Kannada for a boat.
\end{quote}

\subsection*{You'll want to know: \kn{ಕಣ್ಣೀರು}}
\textbf{\kn{ಕಣ್ಣೀರು}} — \emph{kaṇṇīru} — \textquotedblleft{}tears\textquotedblright{}.

\begin{grammarlens}[title={What you've built}]
One more word for this chapter.
\end{grammarlens}

\subsection*{Your turn}
\begin{itemize}
\raggedright
  \item \emph{Say it:} \emph{kaṇṇīru}
  \item \emph{Say it:} \emph{kaṇṇīru}, once more
  \item \emph{Say it:} say \emph{kaṇṇīru}
  \item \emph{Recall:} say the Kannada for wheat, then the Kannada for a boat, then say \emph{kaṇṇīru} again
\end{itemize}

\begin{tcolorbox}[breakable,colback=teal!4,colframe=teal!35!black,title={Before you move on}]
What does \kn{ಕಣ್ಣೀರು} mean? (\textquotedblleft{}Tears\textquotedblright{}.) Say it once more.
\end{tcolorbox}
\section[mīse]{\kn{ಮೀಸೆ} (mīse) — a moustache}

\label{lesson:KA-C139-mise}



\begin{quote}
Before the new one: say the Kannada for a boat, then the Kannada for tears.
\end{quote}

\subsection*{You'll want to know: \kn{ಮೀಸೆ}}
\textbf{\kn{ಮೀಸೆ}} — \emph{mīse} — \textquotedblleft{}a moustache\textquotedblright{}.

\begin{grammarlens}[title={What you've built}]
One more word for this chapter.
\end{grammarlens}

\subsection*{Your turn}
\begin{itemize}
\raggedright
  \item \emph{Say it:} \emph{mīse}
  \item \emph{Say it:} \emph{mīse}, once more
  \item \emph{Say it:} say \emph{mīse}
  \item \emph{Recall:} say the Kannada for a boat, then the Kannada for tears, then say \emph{mīse} again
\end{itemize}

\begin{tcolorbox}[breakable,colback=teal!4,colframe=teal!35!black,title={Before you move on}]
What does \kn{ಮೀಸೆ} mean? (\textquotedblleft{}A moustache\textquotedblright{}.) Say it once more.
\end{tcolorbox}
